\section{应用：无标号有根树与 Pólya 方程}


这一节把 Burnside/Pólya 用到“无标号树”的计数：同构的树算一种。核心做法是把“根的孩子”看成一个\textbf{无序多重集}，用循环指数写生成函数方程。

\textbf{本节用到的记号：} $R(x)$ 为无标号有根树（每个结点至多 $3$ 个无序孩子；含空树）的普通生成函数；$R(x^k)$ 表示把 $x$ 换成 $x^k$；$x$ 记录结点数。

\subsection{为什么不能用 EGF}


带标号树的结点带 $1,\dots,n$ 的标签，$n$ 个结点的结构数用 EGF 很自然；无标号树把同构的树视为同一个，标签没了，\textbf{不能用 EGF 先分标签}。无标号结构要用普通生成函数，并用 Pólya 处理“孩子无序”的对称性。

\subsection{组合方程}


设 $\mathcal R$ 为“无标号有根树”类，每个结点至多有 $3$ 个无序孩子。记

\[
R(x)=\sum_{n\ge0}r_nx^n,
\]

其中 $r_0=1$ 代表\textbf{空树}（用于填充“空的孩子槽”），$r_n$ 为 $n$ 个结点的非空树个数。

\textbf{把每个结点看成“根 + 3 个孩子槽”.} 每个槽可以是空的（贡献 $1$）或一棵树（贡献 $R(x)$），三个槽\textbf{无序}。由 $S_3$ 的循环指数

\[
Z_{S_3}(t_1,t_2,t_3)=\frac16\left(t_1^3+3t_1t_2+2t_3\right),
\]

代入 $t_i=R(x^i)$（槽的轮换长度为 $i$ 时必须放同一棵树，故把 $x$ 换成 $x^i$），再乘根自己的 $x$，得到

\[
R(x)=1+\frac{x}{6}\left(R(x)^3+3R(x)R(x^2)+2R(x^3)\right).
\]

这里最前面的 $1$ 对应空树；右端的三个槽允许为空，所以“至多 $3$ 个孩子”被自动处理。

\textbf{另一种等价写法.} 若把 $R$ 定义为\textbf{非空}树的生成函数（常数项为 $0$），则方程写成

\[
R(x)=x\left(1+R(x)+\frac{R(x)^2+R(x^2)}{2}
+\frac{R(x)^3+3R(x)R(x^2)+2R(x^3)}{6}\right),
\]

两者只差“是否把空树算进 $R$”，系数相同。

\subsection{前几项}


把方程右端反复迭代（或按 $x$ 的幂逐项比较），得

\[
R(x)=1+x+x^2+2x^3+4x^4+8x^5+17x^6+39x^7+89x^8+\cdots
\]

即 $n=1,2,\dots$ 个结点的无标号有根树（每个结点至多 $3$ 个无序孩子）分别有

\[
1,1,2,4,8,17,39,89,\dots
\]

种。这个方程没有简单的闭式，通常用递推/迭代计算系数。

\subsection{与带标号情形的对照}


\begin{itemize}
\item \textbf{带标号、至多 2 个无序孩子}：EGF $T=x(1+T+T^2/2)$，计数 $1,2,9,60,540,\dots$（从 $n=1$ 起），乘法是二项卷积；
\item \textbf{无标号、至多 3 个无序孩子}：OGF 满足上面的 Pólya 方程，计数 $1,1,2,4,8,17,\dots$（从 $n=1$ 起），乘法是普通卷积 + 对称性修正。
\end{itemize}


两者都描述“孩子无序的树”，但一个带标签、一个不带标签，方程与数列完全不同；做题时先问“同构算不算同一种”。

\subsection{树的自同构与染色计数}


若给定一棵具体的树 $T$，想数它用 $C$ 种颜色的染色在 $\operatorname{Aut}(T)$ 作用下的轨道数，也可用 Burnside：

\[
\#(\text{染色轨道})=\frac{1}{|\operatorname{Aut}(T)|}
\sum_{\sigma\in\operatorname{Aut}(T)}C^{c(\sigma)},
\]

其中 $c(\sigma)$ 是 $\sigma$ 作为顶点置换的轮换个数。若还关心各颜色的使用次数，就把 $C^{c(\sigma)}$ 换成 $\left(\sum_j w_j\right)^{c(\sigma)}$ 的循环指数形式。这与 Pólya 公式是同一件事，只是群换成树的对称群。

\subsection{本节的提醒}


\begin{itemize}
\item \textbf{无标号与带标号不同}：不要对无标号树直接套 EGF 的 $\exp$。
\item \textbf{孩子无序时}，$2$ 棵与 $3$ 棵孩子的项要分别用 $S_2,S_3$ 的循环指数，不能只写 $R^2/2$、$R^3/6$（还要加 $R(x^2),R(x^3)$ 的修正项）。
\item \textbf{$R(x^k)$ 的系数不是 0/1 的简单重复}：它表示“同一棵子树重复 $k$ 次”的贡献，展开时按 $x$ 的幂合并。
\item 这类方程一般无闭式；考试若要求前几项，按系数逐项递推即可。
\end{itemize}
